\documentclass[conference]{IEEEtran}
\IEEEoverridecommandlockouts
\usepackage[T1]{fontenc}
\usepackage{cite}
\usepackage{amsmath,amssymb,amsfonts}
\usepackage{algorithmic}
\usepackage{graphicx}
\usepackage{textcomp}
\usepackage{xcolor}
\usepackage{booktabs}
% This TeX Live install ships no Type1 Courier (psnfss's default \texttt
% font); fall back to Computer Modern Typewriter, which is present, rather
% than require a system-wide font package install.
\renewcommand{\ttdefault}{cmtt}
\def\BibTeX{{\rm B\kern-.05em{\sc i\kern-.025em b}\kern-.08em
    T\kern-.1667em\lower.7ex\hbox{E}\kern-.125emX}}
\begin{document}

% TODO before submission: this template is generic IEEEtran/conference -
% confirm the target venue/CFP with Dr. Rana and swap in that conference's
% own style file / page limit if different from the default IEEEtran look.

\title{Diagnosing Domain Mismatch in LLM Security Guardrails:\\
A Deterministic Schema Check for SOC Alert Triage}

\author{
\IEEEauthorblockN{Asma Imran Chaudhry\IEEEauthorrefmark{1}}
\IEEEauthorblockA{
\IEEEauthorrefmark{1}TECIP, Scuola Superiore Sant'Anna \\
CNIT/PNTLab Pisa \\
Pisa, Italy \\
\texttt{asmaimran2003@gmail.com}
}
% TODO: ORCID intentionally left blank -- per project issue #16, still a
% supervisor-level decision deferred until closer to submission.
% TODO: add Dr. Rana Abu Bakar as co-author here if/when issue #16's
% authorship question is settled -- do not assume single-authorship.
}

\maketitle

\begin{abstract}
Security operations centers increasingly rely on large language models (LLMs) to triage the
alert volumes that outpace manual review, but LLM-based triage pipelines inherit LLM-specific
risks such as prompt injection through untrusted alert fields. A natural mitigation is a
learned guardrail that screens alert text before it reaches the model. We report a case study
in which exactly this approach failed silently on real data, why it failed, and what replaced
it. A TF-IDF and logistic-regression classifier trained to detect prompt-injection attempts in
an alert-title field scored 0.883 F1 on a chat-style jailbreak benchmark, then achieved only
52.5\% best-case accuracy when evaluated against real Security Incident Prediction (GUIDE)
data. We show the root cause is not a training-data mismatch in style, as first suspected, but
a schema mismatch: the target field is a numeric identifier in the real dataset, not free text
at all, so no amount of domain-matched fine-tuning could have closed the gap. We replace the
learned classifier with a deterministic type check that separates the same two classes by
construction, verify it against the field's full value space (all 9{,}516{,}837 rows of the
real training data, zero false positives), and report its cost (microsecond-scale) against the
classifier it replaces. We also report an unrelated but related reliability finding: a later
routing change silently disconnected part of the pipeline's graph, producing a hardcoded default
verdict for every alert with no runtime error, caught only by a post-hoc audit -- and describe the
structural regression tests now guarding against a recurrence. We evaluate the full triage agent
against a Random Forest baseline on real GUIDE data and report the result honestly: the
LLM/RF-fallback hybrid does not outperform the RF baseline's raw classification metrics on the
largest sample we ran, current model and formatting-bug fix applied (macro F1 0.648 vs.\ 0.751,
$n=999$), a result we discuss rather than omit.
\end{abstract}

\begin{IEEEkeywords}
security operations center, large language models, prompt injection, guardrails, alert triage,
MITRE ATT\&CK, negative results
\end{IEEEkeywords}

\section{Introduction}

Security operations center (SOC) analysts triage large volumes of alerts under time pressure,
and most alerts turn out to be benign \cite{ndichu2026survey}. LLM-assisted triage -- using a
language model to read an alert, pull in relevant context, and produce a plain-language verdict
with reasoning -- has been shown to reduce analyst workload in both industry deployments
\cite{msseccopilot2024} and recent research systems \cite{cortex2025}. Because these systems put
an LLM in the path of attacker-influenced input (alert titles, device names, and other fields
that may originate from a compromised or adversarial source), they inherit LLM-specific risks:
prompt injection, where crafted text in an alert field attempts to override the model's
instructions rather than describe a genuine security event \cite{ferrag2024genai}.

A standard mitigation is to add a guardrail stage before the LLM sees untrusted text -- either
deterministic pattern matching or a learned classifier. This paper documents a case where we
built and evaluated both, expected the learned classifier to be the stronger of the two, and
found instead that it had no usable signal on real production data. The reason turned out not to
be what we first assumed. We present this as a full account, including the incorrect first
hypothesis, because the eventual root cause -- a schema property of the underlying dataset, not
a text-domain gap -- is the kind of finding that a purely metric-driven evaluation would not have
surfaced on its own.

Our contributions are:
\begin{itemize}
\item A worked example of diagnosing \emph{why} a security classifier fails on real data, going
past accuracy numbers to the actual data schema, and showing the failure mode was structural
(a numeric ID field, not text) rather than a distributional one (informal vs.\ formal language).
\item A deterministic replacement guardrail that separates the same two classes by construction
instead of by learned approximation, verified against the field's entire real value space rather
than a sample.
\item An honest evaluation of the full triage pipeline against a classical baseline on real data,
including a finding -- a silent graph-wiring regression that made every alert default to the same
verdict -- that a purely accuracy-based evaluation had already missed once, and the structural
tests we added to stop it recurring.
\item A candid comparison of the LLM/RF-hybrid triage agent against the Random Forest baseline
that does not outperform it on our largest real-data sample, discussed rather than omitted.
\end{itemize}

\section{Related Work}

Industry SOC copilots such as Microsoft Security Copilot \cite{msseccopilot2024} demonstrate that
LLMs integrated into SOC workflows reduce analyst triage time, but as a closed commercial system
it offers no visibility into guardrail design or failure modes. Ferrag et al.\
\cite{ferrag2024genai} survey LLM applications across the cybersecurity lifecycle and separately
catalog LLM-specific vulnerabilities including prompt injection, which motivates the guardrail
design examined in this paper. Srinivas et al.\ \cite{socaugsurvey2025} similarly survey LLMs and
agentic systems for SOC security automation, situating this paper's single-pipeline design within
that broader landscape. CORTEX \cite{cortex2025} argues for multi-agent, evidence-grounded
triage over a single do-everything model, closely matching our own MITRE ATT\&CK-grounded
retrieval design (Section~\ref{sec:architecture}); we adopt a hybrid LLM/classical-ML routing
strategy rather than multi-agent decomposition, and this paper's focus is the guardrail layer
that sits in front of that routing rather than the routing architecture itself. Ndichu et al.\
\cite{ndichu2026survey} survey AI-driven alert screening broadly and review 22 alert-level
datasets by representational gaps against real SOC environments; their taxonomy is the reference
point we use in Section~\ref{sec:discussion} to discuss the limits of evaluating on a single
static dataset. Ismail et al.\ \cite{wazuh2025rag} build a retrieval-augmented LLM copilot
directly on the open-source Wazuh SIEM's live alert stream -- architecturally the closest system
to ours, but grounded in a different (live, rule-engine-generated) alert source than the static
dataset snapshot used here; we discuss this gap in Section~\ref{sec:futurework}. Red-team-side
agentic tooling such as ADStrike \cite{adstrike2025} demonstrates the offensive counterpart of
the agentic-AI pattern this paper applies defensively. GeNIS \cite{genis2025} addresses a
labeled-data gap for small/medium-enterprise network traffic that our own dataset (Section
\ref{sec:dataset}) does not cover, and is discussed as future work.

\section{System Architecture}
\label{sec:architecture}

The triage pipeline is implemented as a directed graph (LangGraph) over a shared alert state.
Given a raw alert, the pipeline: (1) builds a structured text context from the alert's populated
fields; (2) applies a deterministic regex-based input guardrail; (3) applies the schema guardrail
described in Section~\ref{sec:guardrail}; (4) retrieves MITRE ATT\&CK technique context for any
technique ID present on the alert; (5) routes the alert to either the LLM classifier or a
Random-Forest fallback classifier, depending on how much analyst-readable evidence the alert
carries; and (6) parses the resulting verdict, routing low-confidence results to a human-review
checkpoint rather than auto-closing them. A failure or block at any guardrail stage routes
directly to human review rather than reaching the LLM.

The routing decision at step 5 exists because a meaningful fraction of real GUIDE alerts carry
almost no analyst-readable evidence: MITRE technique, suspicion level, and last-verdict fields
are frequently absent, and \texttt{AlertTitle}/\texttt{Category} alone are numeric or
closed-vocabulary codes rather than descriptive text (this is the same schema property discussed
in Section~\ref{sec:rootcause}). Alerts with fewer than two of these three discriminative fields
populated are routed to the RF classifier instead of the LLM, on the reasoning that an LLM has
little basis for a per-alert judgment when the alert itself carries almost no readable signal.

\section{The Domain-Mismatch Investigation}
\label{sec:guardrail}

\subsection{First guardrail: a learned classifier}

The regex-based input guardrail (Section~\ref{sec:architecture}) catches literal
instruction-override and role-impersonation phrasing, but by construction cannot catch paraphrased
or novel injection attempts. As a second-stage defense, we trained a TF-IDF + logistic-regression
classifier on \texttt{AlertTitle} text, selected after benchmarking against Meta's Prompt Guard and
Protect AI's LLM Guard on a chat-style jailbreak benchmark, where it scored 0.883 F1 -- the
strongest result of the three.

\subsection{Real-data evaluation and a masking bug}

Evaluated against real GUIDE alert data rather than the benchmark it was selected on, the
classifier initially flagged 100\% of real alerts as injection attempts. Investigating, we found
an index bug in how the classifier's output probability was read (\texttt{predict\_proba(X)[0][0]},
the benign-class probability, was being used where \texttt{[0][1]}, the injection-class
probability, was needed) -- caught in code review before the initial finding was accepted at face
value. After the fix, benign and injection scores on real alert text were statistically
indistinguishable (mean 0.209 vs.\ 0.224); the best achievable accuracy across a threshold sweep
was 52.5\%, at chance level for a binary decision.

\subsection{Root cause}
\label{sec:rootcause}

The first hypothesis was a training-data style mismatch: the classifier had been benchmarked on
conversational jailbreak text, and real SOC alert text might simply look different enough
linguistically that the model's learned features did not transfer. Before spending the following
week's time-boxed effort on a domain-matched fine-tune, we checked what \texttt{AlertTitle}
actually contains in the real data. It is not text at all: across all 9{,}516{,}837 rows of the
real training set, \texttt{AlertTitle} is a numeric ID field with 86{,}149 unique integer values
(we re-verified this figure directly against the raw CSV during a later audit pass, confirming it
independently rather than propagating it from the original finding). GUIDE's own published
description of the alert dataframe lists \texttt{OrganizationId}, \texttt{DetectorId},
\texttt{ProductId}, \texttt{Category}, and \texttt{Severity} as its only categorical columns --
no free-text field exists in the schema at all. A TF-IDF classifier cannot find linguistic signal
in a field that was never linguistic; no amount of domain-matched training data would have closed
that gap, because the gap was never about domain.

\subsection{Deterministic schema guardrail}

Since \texttt{AlertTitle} and \texttt{DetectorId} are supposed to always be numeric IDs, the
correct second-stage check is not a classifier but direct type validation: does the field parse
as an integer. An injection payload is, by definition, not a valid integer, so this check
separates the two classes by construction rather than by learned approximation, and costs
microseconds rather than a model load. We tested it two ways: on a synthetic 20-benign/20-injection
set (100\% accuracy, matching the earlier classifier's benchmark-set performance but now on the
real risk case rather than a proxy for it), and, during the audit for this paper, directly against
the field's entire real value space -- all 9{,}516{,}837 \texttt{AlertTitle} values in the training
set -- confirming zero non-numeric values and therefore zero false positives across the complete
population, not a sample of it. Unlike the learned classifier, this check was hard-gated into the
pipeline: it has no threshold, no accuracy trade-off, and no probabilistic gray zone to be
cautious about.

\section{Evaluation}

\subsection{Dataset}
\label{sec:dataset}

We use the Microsoft GUIDE dataset (CDLA-2.0) \cite{guide2024}, a public release of 1.6M real SOC alerts and 1M
analyst-triaged incidents from over 6{,}100 organizations. All quantitative results in this paper
were produced against the real dataset, not the schema-matched synthetic sample used for local
development, and are reproducible from the committed artifacts in the project repository.

Published alert-triage evaluations on large real-world datasets commonly bound their evaluation to
a subsample well below full dataset scale, though the reason for bounding differs by evaluation
type. Freitas et al., who introduced GUIDE itself, restrict their evaluation of the investigation
module to 1{,}000 randomly selected incidents out of GUIDE's roughly one million -- not an
LLM-based system, but a classical (Random Forest, PCA, and cosine-similarity) approach whose
evaluation required manual analyst relevance judgments per incident, which bounds feasible scale
independent of inference cost -- despite classical statistical baselines in the same paper being
evaluated on full regional splits an order of magnitude larger \cite{freitas2024guide}. Zhao et
al.\ evaluate LLM-based alert triage, where per-call inference cost and rate limits are the binding
constraint, on stratified test splits of roughly 785 and 20{,}000 alerts drawn from two
production-derived corpora \cite{zhao2025infodense}. Our own evaluation samples of 300 and up to
999 alerts (Section~\ref{sec:hybrideval}), likewise bounded primarily by LLM inference cost and
rate limits, are consistent in order of magnitude with Zhao et al.'s LLM-specific evaluation scale,
and, on the same source dataset, close to the 1{,}000-incident scale GUIDE's own authors used for
their (non-LLM) investigation-module evaluation -- while remaining below the larger,
tens-of-thousands-scale subsamples some papers draw from smaller base corpora. We report this
comparison for calibration, not as a claim that a
larger sample would be unnecessary; Section~\ref{sec:futurework} discusses scaling the evaluation
further as the API-quota constraint that currently bounds it (Section~\ref{sec:redteam}) allows.

\subsection{RF baseline}

A Random Forest (200 trees) trained on 79{,}580 alerts (stratified 80/20 split, $n=100{,}000$,
fixed random seed) and evaluated on a 19{,}895-alert holdout achieves macro F1 0.751 (accuracy
0.772). We re-trained this baseline from scratch during the audit for this paper (rather than
trusting the previously saved artifact) and reproduced macro F1 0.751 exactly.

\subsection{Hybrid agent vs.\ RF baseline}
\label{sec:hybrideval}

On real GUIDE alerts, the full LLM/RF-hybrid pipeline was evaluated at two sample sizes. On a
300-alert sample, accuracy was 0.753 and macro F1 0.750 (81.3\% of alerts routed to the RF
fallback path, 18.7\% to the LLM), closely matching the RF baseline. On a larger, later 999-alert
sample -- the largest real-data evaluation in the project, and the one we treat as primary here --
accuracy was 0.646 and macro F1 0.648 (79.1\% RF-routed, 20.9\% LLM-routed), a roughly 10-point
macro F1 gap below the RF baseline alone. This number reflects the current model
(\texttt{openai/gpt-oss-20b} via Groq) and the
\texttt{build\_context()} formatting-bug fix (Section~\ref{sec:llmgap}); an earlier version of this
same 999-alert run, on the since-retired \texttt{llama-3.1-8b-instant} model and with the
formatting bug present, scored macro F1 0.669 -- Section~\ref{sec:llmgap} discusses why neither
fix closed the gap, and if anything made it marginally wider. We report both sample sizes rather
than the more favorable of the two: the smaller sample was the one originally written up, and the
larger, more statistically robust sample was sitting in the repository's committed results without
ever having been discussed. We do not read this as a failure of the hybrid design so much as a
reminder that raw
classification metrics are not the only axis this system is designed against -- the LLM path adds
per-alert natural-language reasoning and MITRE-grounded context that the RF baseline does not
produce at all, which the macro F1 number does not capture, but which is exactly the
analyst-facing output the system exists to produce. We do not have an analyst-rated evaluation of
that reasoning quality, and flag its absence as a limitation (Section~\ref{sec:discussion}) rather
than assert the trade-off is worth it without evidence.

\subsection{Where the hybrid gap comes from, and whether it closes}
\label{sec:llmgap}

Splitting an earlier version of the 999-alert run (deprecated model, formatting bug present) by
routing path isolates the source of the gap above: the RF-routed subset (790 alerts, 79.1\%) alone
scores macro F1 0.752, matching the standalone RF baseline almost exactly, while the LLM-routed
subset (209 alerts, 20.9\%) scores macro F1 0.308 -- barely above chance for a three-class problem.
This is the entire gap; RF's own accuracy is not the problem. It
is also a counterintuitive result on its face, since \texttt{route\_by\_context} sends the LLM
specifically the alerts with the \emph{most} analyst-readable context (MITRE technique, suspicion
level, prior verdict; Section~\ref{sec:architecture}) -- the alerts an LLM should be best suited
for, not worst.

We investigated two candidate explanations with live re-evaluation on this same 209-alert subset,
both during this audit. First, that 999-alert run used \texttt{llama-3.1-8b-instant}, since retired
from Groq's catalog (Section~\ref{sec:redteam} discusses its unrelated discovery); a stronger
current model might close the gap on its own. Second, inspecting \texttt{build\_context()} found a
real formatting bug: a bare truthiness check (\texttt{if alert.get(field):}) does not treat a
missing-data \texttt{NaN} float as absent, so a missing MITRE technique or suspicion level was not
omitted from the prompt but rendered as the literal text ``MITRE Technique: nan'' -- present in 55\%
and 31\% of this subset's prompts respectively. We fixed the formatting bug unconditionally (it has
no plausible upside) and then re-ran the LLM-routed subset (60 of the 209 alerts, stratified) with
the current model under the original prompt: accuracy 0.283, macro F1 0.151 -- neither issue was
the cause, and the corrected baseline is, if anything, worse than the deprecated-model number on the
same alerts. Reading the model's own reasoning text confirms why: it is not misreading a garbled
field, it is explicitly concluding ``insufficient evidence $\rightarrow$ FalsePositive'' on alerts
whose only evidence is a generic category label and a bare suspicion flag -- a defensible
inference in isolation that happens to collapse onto one class 55 of 60 times, matching the exact
failure mode Section~\ref{sec:architecture}'s low-context RF fallback already exists to route around
for sparser alerts.

We then confirmed this was not a small-sample artifact by re-running the entire 999-alert
evaluation with the current model and the formatting-bug fix applied under the original prompt --
the same run reported as this paper's primary hybrid headline in
Section~\ref{sec:hybrideval} (macro F1 0.648 overall, $n{=}999$). Its LLM-routed subset now
covers the full 209 alerts rather than the 60-alert stratified sample above: accuracy 0.230,
macro F1 0.174 -- consistent with, and slightly worse than, the 60-alert diagnostic's 0.151, and
still well below the deprecated-model number it was meant to explain. The RF-routed subset is
unaffected by either fix and reproduces 0.752, as expected.

We then time-boxed one prompt-engineering attempt on the same 60 alerts: explicit guidance on which
fields carry the strongest signal, three grounded few-shot examples using GUIDE's own field
conventions, and reordering the response schema so reasoning is written before the verdict rather
than after. This raised accuracy to 0.333 and macro F1 to 0.268 -- a real improvement over the
corrected baseline, and it eliminated the single-class collapse (precision/recall are no longer
near-zero on two of three classes) -- but it remains well below RF's 0.752 on the same kind of
alert, and slightly below even the deprecated-model number this investigation set out to explain.
We read this as a genuine, if partial, negative result rather than a tuning failure: a
few-hundred-token prompt cannot reliably supply what GUIDE's schema does not encode for these
alerts, and further prompt iteration is unlikely to be the more productive investment relative to
richer context retrieval (Section~\ref{sec:futurework}). We report the 0.751-vs-0.152-vs-0.268
progression in full rather than presenting only the improved number, since the honest conclusion --
that neither a newer model nor a stronger prompt closes this specific gap -- is itself evidence
relevant to Section~\ref{sec:discussion}'s discussion of when a hybrid LLM/RF design is and is not
justified on classification accuracy alone.

\subsection{Guardrail cost and accuracy}

The regex input guardrail runs in a mean 3.6~\textmu s per check (10{,}000-iteration
microbenchmark), blocking 10{,}000/10{,}000 known injection strings with zero false positives on
benign alerts in the same benchmark. The schema guardrail (Section~\ref{sec:guardrail}) adds a
second, equally cheap deterministic stage in front of the LLM/RF routing decision.

\subsection{Throughput}

On real-alert throughput benchmarks, the LLM path sustains roughly 0.37--0.57 alerts/second (mean
latency 1.8--2.7~s per alert), dominated by external API latency and largely unaffected by added
worker concurrency because the bottleneck is the remote provider, not local compute. The RF path,
by contrast, sustains 16.7--41.9 alerts/second at the same prompt volume, scaling with local
worker count as expected for CPU-bound inference. This is the practical justification for routing
low-evidence alerts to RF rather than LLM: those alerts gain little from LLM reasoning
(Section~\ref{sec:architecture}) and the RF path is roughly two orders of magnitude faster.

\subsection{Reliability engineering: a silent routing regression}

A later change that introduced the schema guardrail into the pipeline graph also, unintentionally,
deleted the conditional edge routing MITRE-enriched alerts onward to the LLM/RF classification
step, without replacing it. The graph still compiled -- the underlying graph library does not
validate that every node has an outgoing edge -- and ran to completion on every alert, but produced
no verdict at all past that point. The evaluation harness's own code, reading the missing verdict
field with a default value, silently treated every alert as a low-confidence \texttt{FalsePositive}
guess instead of raising an error. No evaluation had been re-run since the change landed, so
nothing surfaced the regression until a manual, unrelated documentation audit found it days later.
We root-caused it against the exact commit that introduced it, restored the missing edge, and
added a structural regression test suite (the project's first) that asserts every non-terminal
node in the pipeline graph has a reachable outgoing edge, plus an end-to-end invocation test
asserting a real verdict is produced. We consider this worth reporting alongside the classifier
finding: both are cases where a plausible-looking result (a classifier's benchmark score in one
case, a batch of predictions in the other) was masking a structural problem that only became
visible on inspecting the underlying mechanism rather than the output metric.

\subsection{Adversarial robustness: a dynamic red-team evaluation}
\label{sec:redteam}

The evaluations above measure classification accuracy against labeled data, not whether an
attacker-controlled alert field can manipulate the triage verdict itself. To test that question
directly, we used the open-source deepteam framework \cite{deepteam2025} to dynamically generate
adversarial inputs against \texttt{classify\_with\_llm} in isolation, bypassing both guardrails
deliberately -- they already have dedicated passing tests, and isolating the LLM node answers the
narrower question of whether the model itself is manipulable once adversarial content reaches it.
Since deepteam has no built-in Groq provider, we implemented a Groq-backed judge/simulator wrapper,
used as both roles. Three vulnerability types (hijacking, social-engineering goal theft, and
cross-context indirect instruction) were crossed with four attack transforms (prompt injection,
Base64, ROT13, and roleplay), yielding 12 test cases.

Of 12 cases, 7 completed and were conclusively scored, with a 0\% attack success rate -- every
completed attack was resisted. The remaining 5 errored, concentrated entirely in the two attack
methods (prompt injection, roleplay) whose generation step itself requires a structured JSON
response from the judge model; the two deterministic transforms completed 3/3 each with zero
errors. We traced the cause to the judge model producing malformed JSON and added a defensive
JSON-repair layer to the wrapper. A same-day retry was blocked by the judge model's daily token
quota being exhausted from cumulative testing, leaving the fix unit-tested but not live-verified at
the time.

We later live-verified the fix once the quota reset. A focused retry on the two problematic attack
methods raised their conclusive rate from a combined 1/6 to 4/12 relative to the original run
(prompt injection 0/3~$\to$~1/6, roleplay 1/3~$\to$~3/6), with all four newly-conclusive cases again
resisting the attack. The fix is a real, measured improvement rather than a full resolution: the
remaining cases still error on the same JSON-formatting failure mode, indicating the judge model's
structured-output reliability is a persistent limitation of this evaluation setup rather than a bug
we introduced.

We also ran deepteam against the full defended pipeline, guardrails included, to test whether an
attack that reaches the LLM is still caught downstream. This produced a counterintuitive result on
inspection: 9 of 12 cases were conclusive with 0\% attack success, but every conclusive case had
actually been routed to the RF fallback classifier, not the LLM, because the synthesized test alert
lacked the context fields (MITRE technique, suspicion level, prior verdict) that the pipeline's own
sparse-context routing logic (Section~\ref{sec:architecture}) checks before choosing between the
LLM and RF paths. The guardrails passed cleanly on every case; they simply were not the deciding
factor, since the alert rarely reached the node they exist to protect. We report this as a
test-harness limitation surfaced by inspecting per-case routing rather than trusting the summary
pass rate -- the same lesson as the routing regression above, that an aggregate metric in this
codebase is not sufficient on its own to confirm a mechanism worked as intended.

We fixed this test-harness gap and re-ran the identical scope: the synthesized alert now carries
two deliberately neutral evidence values (a placeholder suspicion level and prior verdict, chosen
to satisfy the routing threshold without leaning the verdict either way), which is enough to clear
the two-of-three threshold and send the alert down the LLM path instead of RF. The fix worked as
intended -- 8 of the 12 cases now show the LLM path in their output, versus 0 of 12 before. Both
guardrails again passed cleanly on every case, confirming they are genuinely not bypassed rather
than merely untested. Of the 8 cases that reached the LLM, 6 completed and were scored, and all 6
resisted the attack (0\% attack success on conclusive cases), consistent with the LLM-only result
above. The remaining 6 errors reduce, on inspection, to the same judge-model JSON-reliability
limitation already discussed, not a new failure mode introduced by the fix.

\section{Discussion and Limitations}
\label{sec:discussion}

\textbf{Is the LLM path worth it?} On classification accuracy alone, our evidence says no: the
RF-routed subset of every hybrid run we have matches the standalone RF baseline, while the
LLM-routed subset never does, and neither a newer model nor a time-boxed prompt-engineering pass
closed that gap (Section~\ref{sec:llmgap}). If macro F1 on GUIDE's three-way label were the only
criterion, the honest recommendation from this evidence is the RF baseline alone, not a hybrid. What
a hybrid design adds instead is qualitative and mostly untested by us directly: MITRE-grounded,
per-alert natural-language reasoning that a Random Forest cannot produce at all, and -- based on the
full-graph red-team finding (Section~\ref{sec:redteam}) -- distinct behavior on alerts with little
structured evidence, where the RF path is not exercised and the LLM path must reason from sparse
text rather than fail closed. We did not run an analyst-rated evaluation of that reasoning's
usefulness, so we do not assert the trade-off is worth it; we report the accuracy cost precisely so
that judgment can be made on real numbers rather than an assumed LLM advantage. A defensible
narrower claim survives this evidence: routing by evidence density, as this pipeline already does,
is the correct design regardless of the LLM subset's weak accuracy, since it already sends the
LLM only the alerts RF is not confidently exercised on, and confines the accuracy cost to exactly
that subset rather than the whole pipeline.

The schema guardrail's guarantee is narrower than it may appear: it verifies a field is
\emph{an} integer, not that it is a \emph{plausible} one. A crafted negative or out-of-range
integer passes the check even though no real alert in the dataset has ever presented one,
because the check's only job -- separating free text from IDs -- does not require checking
plausibility, and injection payloads are categorically excluded by the integer check regardless.
We note this as a documented scope boundary rather than an oversight: a different guardrail would
be needed to catch a numeric-payload attack that this check was never designed to catch.

The RF baseline and fallback classifier were trained exclusively on GUIDE's feature schema. We
built (as part of this same research program, reported alongside this paper) a schema adapter
translating alerts from Wazuh, an open-source SIEM, into the pipeline's alert format, to test
generalization beyond one static dataset snapshot. The LLM/guardrail path handles Wazuh-derived
alerts correctly -- confirmed by direct invocation, not merely a unit test asserting the mapping
function's output shape. The RF fallback path does not fail outright on Wazuh-derived alerts (the
underlying classifier's missing-value handling accepts the input and returns a normal-looking
prediction), but a large majority of its expected feature columns are absent for a Wazuh-sourced
alert -- a sparsity pattern the model was never trained or validated against, unlike GUIDE's own
in-distribution missingness. We flag this explicitly as an open validity question rather than
either fixing it speculatively or hiding it: routing a Wazuh-derived alert through the RF fallback
currently produces a result of unknown reliability, and we do not report an accuracy number for it
because we do not have labeled Wazuh data to check one against.

Evaluating against a single dataset snapshot -- even a large, real one -- limits how confidently
any of these results generalize; Ndichu et al.\ \cite{ndichu2026survey}'s survey of 22 alert-level
datasets is instructive here in cataloging just how differently existing benchmarks tend to be
structured relative to live SOC environments. GeNIS \cite{genis2025}, a 2025 dataset built
specifically to address a labeled-data gap for small/medium-enterprise network traffic, is a
candidate for extending this evaluation beyond GUIDE, but its flow-level schema does not map onto
GUIDE's incident/alert schema and would require its own loader rather than an extension of the
existing one -- noted as a scoping decision for future work rather than attempted here.

\section{Future Work}
\label{sec:futurework}

Three extensions follow directly from the limitations above. First, a live Wazuh deployment
(rather than the schema-only adapter reported here) would let the pipeline's guardrails and
routing be tested against a continuously generated alert stream instead of a static snapshot,
closer to the deployment described by Ismail et al.\ \cite{wazuh2025rag}. Second, integrating
GeNIS \cite{genis2025} as a second evaluation dataset -- once its flow-level schema is reconciled
with an incident/alert-level loader -- would test whether the RF baseline and schema guardrail's
strong GUIDE-specific results hold on a structurally different, SME-focused dataset; this remains
an explicitly deferred scoping decision (Section~\ref{sec:discussion}) rather than an in-progress
item, pending supervisor sign-off on the schema-reconciliation approach. Third, the RF-fallback
reliability gap on non-GUIDE alert sources (Section~\ref{sec:discussion}) motivates either labeled
data collection on Wazuh-origin alerts specifically, or an origin-aware routing change that
defaults non-GUIDE alerts to human review rather than the RF path until validated. The full-graph
red-team harness gap noted in an earlier draft of this section -- the synthesized alert not
carrying a context field \texttt{route\_by\_context} checks -- has since been fixed and re-run
(Section~\ref{sec:redteam}); the remaining unresolved judge-model JSON-reliability errors are a
still-open candidate for a stronger or differently-hosted judge model if one becomes available.

\section{Conclusion}

We diagnosed why a learned prompt-injection classifier had no usable signal on real SOC alert
data: not a domain mismatch in text style, as first suspected, but a schema mismatch -- the target
field is a numeric ID, not text. We replaced it with a deterministic check verified against the
field's complete real value space, at a fraction of the computational cost. We report this
alongside an unrelated reliability finding from the same investigation -- a silent routing
regression masked by an evaluation harness's own default-value handling -- because both point to
the same lesson: inspecting the underlying data and control flow directly, rather than trusting an
output metric, is what surfaced each problem. We evaluate the resulting pipeline against a
classical baseline honestly, including a real-data comparison where the hybrid pipeline does not
outperform the baseline on raw classification metrics, and discuss what that trade-off does and
does not tell us.

\section*{Acknowledgment}
% TODO: fill in per your institution's / supervisor's guidance -- deferred
% per issue #16 pending funding/acknowledgment decisions.

\begin{thebibliography}{00}

\bibitem{msseccopilot2024} Microsoft, ``Microsoft Security Copilot,'' Microsoft, 2024. [Online].
Available: https://www.microsoft.com/en-us/security/business/ai-machine-learning/microsoft-copilot-for-security

\bibitem{adstrike2025} capture0x, ``ADStrike: AI-Powered Modular Active Directory Red-Team
Framework via MCP,'' GitHub repository, 2025. [Online]. Available:
https://github.com/capture0x/AdStrike

\bibitem{ferrag2024genai} M. A. Ferrag, F. Alwahedi, A. Battah, B. Cherif, A. Mechri, N. Tihanyi,
T. Bisztray, and M. Debbah, ``Generative AI in Cybersecurity: A Comprehensive Review of LLM
Applications and Vulnerabilities,'' arXiv:2405.12750, 2024.

\bibitem{cortex2025} B. Wei, Y. S. Tay, H. Liu, J. Pan, K. Luo, Z. Zhu, and C. Jordan, ``CORTEX:
Collaborative LLM Agents for High-Stakes Alert Triage,'' arXiv:2510.00311, 2025.

\bibitem{socaugsurvey2025} Srinivas, Kirk, Zendejas, Espino, Boskovich, Bari, Dajani, and
Alzahrani, ``AI-Augmented SOC: A Survey of LLMs and Agents for Security Automation,'' \emph{J.
Cybersecur. Priv.}, vol. 5, no. 4, Art. 95, 2025.

\bibitem{wazuh2025rag} Ismail, R. Kurnia, F. Widyatama, I. M. Wibawa, Z. A. Brata, Ukasyah, G. A.
Nelistiani, and H. Kim, ``Enhancing Security Operations Center: Wazuh Security Event Response with
Retrieval-Augmented-Generation-Driven Copilot,'' \emph{Sensors}, vol. 25, no. 3, Art. 870, 2025.

\bibitem{genis2025} M. Silva, D. Pinto, J. Vitorino, J. Gon\c{c}alves, E. Maia, and I. Pra\c{c}a,
``GeNIS: A Modular Dataset for Network Intrusion Detection and Classification,'' \emph{Data in
Brief}, 2025, doi:10.1016/j.dib.2025.111487.

\bibitem{ndichu2026survey} S. Ndichu, T. Ban, S. Ozawa, T. Takahashi, and D. Inoue, ``AI-Driven
Security Alert Screening and Alert Fatigue Mitigation in Security Operations Centers: A
Survey,'' arXiv:2605.08316, 2026.

\bibitem{guide2024} Microsoft, ``GUIDE: Microsoft Security Incident Prediction Dataset,'' Kaggle,
2024. [Online]. Available:
https://www.kaggle.com/datasets/Microsoft/microsoft-security-incident-prediction

\bibitem{deepteam2025} Confident AI, ``DeepTeam: The Open-Source LLM Red Teaming Framework,''
GitHub repository, 2025. [Online]. Available: https://github.com/confident-ai/deepteam

\bibitem{freitas2024guide} S. Freitas, J. Kalajdjieski, A. Gharib, and R. McCann, ``AI-Driven
Guided Response for Security Operation Centers with Microsoft Copilot for Security,''
arXiv:2407.09017, 2024.

\bibitem{zhao2025infodense} G. Zhao, Y. Zhang, C. Tian, D. Xie, H. Liu, and B. Wang,
``Information-Dense Reasoning for Efficient and Auditable Security Alert Triage,''
arXiv:2512.08169, 2025.

\end{thebibliography}

\end{document}
